% =====================================================================
%  Lecture 01 -- GSM, GPRS and EDGE                      (English edition)
%  Companion: lecture-01-fa.tex.  Build: make lecture LECTURE=01
% =====================================================================
\documentclass[11pt,a4paper]{article}
\providecommand\ACNroot{../..}
\input{\ACNroot/templates/lecture-en}

\ACNset{number}{01}
\ACNset{title}{GSM, GPRS and EDGE}
\ACNset{shorttitle}{GSM, GPRS and EDGE}
\ACNset{titleother}{جی‌اس‌ام، جی‌پی‌آراس و اج}
\ACNset{part}{Part I -- 2G/3G foundations (Lectures 1--4)}
\ACNset{time}{6--8 hours (3 h reading, 3--5 h exercises)}
\ACNset{level}{Introductory to intermediate}
\ACNset{prereq}{Decibels, basic digital modulation, TDMA/FDMA, the OSI model}
\ACNset{gen}{2}

\begin{document}
\ACNcover

% ---------------------------------------------------------------------
\section{Learning objectives}\label{sec:objectives}
After studying this lecture you should be able to:
\begin{objectives}
  \item \textbf{Draw} the GSM reference architecture, name every network element
        (MS, BTS, BSC, MSC, VLR, GMSC, HLR, AuC, EIR) and state the interface
        between each pair (Um, Abis, A, C, D, E, F).
  \item \textbf{Derive} the GSM radio timing (bit period, burst, TDMA frame and
        26-multiframe) from the 270.833\,kbit/s gross rate and \textbf{compute}
        the 22.8\,kbit/s gross rate of a full-rate traffic channel.
  \item \textbf{Explain} the GSM challenge--response authentication with the
        triplet $(\text{RAND},\text{SRES},K_c)$ and identify its security weaknesses.
  \item \textbf{Describe} the GPRS overlay (PCU, SGSN, GGSN, Gb, Gn, Gi, Gr) and the
        PDP context, and \textbf{trace} a PDP context activation step by step.
  \item \textbf{Calculate} per-user GPRS/EDGE throughput from the coding scheme and the
        number of timeslots, and \textbf{evaluate} the trade-off between robustness and rate.
  \item \textbf{Relate} cluster size, path-loss exponent and co-channel interference
        through the first-tier $C/I$ model.
\end{objectives}

% ---------------------------------------------------------------------
\section{Prerequisite knowledge}\label{sec:prereq}
The lecture assumes the decibel arithmetic of the prerequisites review
($x_{\mathrm{dB}} = 10\log_{10}x$), the notion of a modulation symbol and of
bandwidth, and the difference between frequency-division (FDMA) and time-division
(TDMA) multiple access. If you are unsure about any of these, revise them before
continuing; every quantitative result below is built from them.

\paragraph{Key terms} The table lists the vocabulary used throughout the lecture with
its Persian equivalent, so that both editions can be read side by side.
\begin{keyterms}
Mobile Station (MS) & \fa{ایستگاه موبایل} & Handset = mobile equipment (ME) + SIM card. \\
Base Transceiver Station (BTS) & \fa{ایستگاه فرستنده‌گیرندهٔ پایه} & Radio site: transmitters, receivers and antennas of one or more cells. \\
Base Station Controller (BSC) & \fa{کنترل‌گر ایستگاه پایه} & Manages radio resources and handovers of many BTSs. \\
Mobile Switching Centre (MSC) & \fa{مرکز سوئیچ موبایل} & Circuit switch for calls; also runs mobility signalling. \\
Home Location Register (HLR) & \fa{ثبات مکان خانگی} & Permanent subscriber database of the home operator. \\
Visitor Location Register (VLR) & \fa{ثبات مکان مهمان} & Temporary copy of data for subscribers in an MSC area. \\
Authentication Centre (AuC) & \fa{مرکز احراز هویت} & Stores the secret key $K_i$ and generates triplets. \\
Timeslot (TS) & \fa{شیار زمانی} & One of the eight TDMA positions on a carrier. \\
Burst & \fa{بُرست} & The radio waveform transmitted in one timeslot. \\
PDP context & \fa{زمینهٔ PDP} & Packet session state: address, APN, QoS, tunnel IDs. \\
SGSN / GGSN & \fa{گره‌های پشتیبان GPRS} & Serving and gateway packet nodes of the GPRS core. \\
Coding scheme (CS/MCS) & \fa{طرح کدگذاری} & Combination of modulation and channel-code rate. \\
\end{keyterms}

% ---------------------------------------------------------------------
\section{Introduction and motivation}\label{sec:intro}
In the 1980s Western Europe operated several incompatible analogue cellular systems.
A handset bought in one country was useless across the border, every system needed
its own equipment industry, and analogue FM voice offered no protection against
eavesdropping. The engineering problem that GSM\footnote{Originally \emph{Groupe
Spécial Mobile}, later \emph{Global System for Mobile communications}.} set out to
solve was therefore not only a radio problem; it was a \emph{system} problem:
\begin{itemize}
  \item \textbf{Interoperability.} Standardised, open interfaces would let an operator
        buy base stations from one vendor and switches from another.
  \item \textbf{Roaming.} A subscriber identity independent of the handset (the SIM)
        and a home database reachable from any visited network.
  \item \textbf{Spectral efficiency and capacity.} Digital speech coding, channel
        coding and TDMA let several users share one carrier and tolerate more
        interference, so frequencies can be reused more tightly.
  \item \textbf{Security.} Authentication of the subscriber and encryption of the radio link.
\end{itemize}
GSM was designed for circuit-switched voice. When data demand grew in the 1990s,
the same radio network was extended with a packet-switched overlay: \emph{GPRS}
(General Packet Radio Service), and later \emph{EDGE} (Enhanced Data rates for GSM
Evolution), which raised the bit rate per timeslot with a higher-order modulation.
Studying these three steps in one lecture shows a pattern that recurs throughout the
course: \textbf{each generation reuses what it can and adds the minimum new
machinery needed for the new problem.}

% ---------------------------------------------------------------------
\section{Fundamental concepts}\label{sec:fundamentals}

\subsection{Cells, frequency reuse and multiple access}
A mobile network divides its service area into \emph{cells}, each served by one
set of antennas. Because spectrum is scarce, the same carrier frequencies are
\emph{reused} in cells that are far enough apart. GSM combines two multiple-access
methods:
\begin{itemize}
  \item \textbf{FDMA:} the band is divided into carriers spaced \qty{200}{\kilo\hertz}
        apart. In the primary GSM-900 band, the uplink occupies \qtyrange{890}{915}{\mega\hertz}
        and the downlink \qtyrange{935}{960}{\mega\hertz}, giving 124 usable carrier pairs
        with a duplex spacing of \qty{45}{\mega\hertz} \cite{ts45002}.
  \item \textbf{TDMA:} each carrier is divided in time into frames of eight
        timeslots; a full-rate voice user occupies one timeslot in every frame.
\end{itemize}

\begin{definitionbox}{Physical and logical channel}
A \emph{physical channel} is a timeslot number on a carrier (or a hopping sequence of
carriers). A \emph{logical channel} is a type of information --- traffic (TCH) or
control (BCCH, CCCH, SDCCH, SACCH, FACCH) --- that is mapped onto a physical channel
according to a fixed multiframe structure.
\end{definitionbox}

\subsection{Separating identity from equipment}
GSM separates the \emph{subscriber} from the \emph{device}. The subscriber is
identified by the IMSI (International Mobile Subscriber Identity: mobile country
code, mobile network code and a subscriber number, at most 15 digits) stored on the
SIM together with a 128-bit secret key $K_i$. The device is identified by its IMEI.
To avoid sending the IMSI over the air, the VLR assigns a temporary identity, the
TMSI, which is used instead whenever possible \cite{ts23003}.

\subsection{Circuit switching versus packet switching}
Voice in GSM is \emph{circuit switched}: when a call is set up, a timeslot on the
radio, a channel on the Abis and A links and a path through the MSC are reserved
for the whole call, whether or not anybody is speaking. Data in GPRS is
\emph{packet switched}: radio blocks are assigned only when the user has data to
send, and many users share the same timeslots. This difference is the reason GPRS
needs new nodes rather than just new software in the MSC, and it is analysed in
depth in Lecture~2.

% ---------------------------------------------------------------------
\section{Detailed technical explanation}\label{sec:technical}

\subsection{The GSM network elements}
\paragraph{Mobile Station (MS)} The MS is the mobile equipment plus the SIM. The SIM
runs the authentication algorithm A3 and the key-generation algorithm A8 so that
$K_i$ never leaves the card.

\paragraph{Base Transceiver Station (BTS)} The BTS contains the radio transceivers
(TRX) of one or more cells. It performs modulation (GMSK), channel coding,
interleaving, encryption with A5 and the radio measurements that feed handover
decisions. It is intentionally simple: decisions are taken by the BSC.

\paragraph{Base Station Controller (BSC)} One BSC controls tens to hundreds of BTSs.
It allocates radio channels, manages power control and decides most handovers
between its own cells. Together the BTSs and their BSC form the \emph{Base Station
Subsystem} (BSS). Speech transcoding between the 13\,kbit/s radio codec and the
64\,kbit/s PCM format of the core is performed in a transcoder unit (TRAU), usually
located at the MSC site to save transmission capacity.

\paragraph{Mobile Switching Centre (MSC) and VLR} The MSC is a telephone exchange
extended with mobility functions: it sets up and releases calls, controls inter-BSC
handovers, and runs location updating and authentication. The VLR (in practice
integrated in the MSC) caches the subscription data and security triplets of every
subscriber currently in its area.

\paragraph{Gateway MSC (GMSC)} Calls arriving from the fixed network (PSTN/ISDN) are
routed to a GMSC, which interrogates the HLR to find where the called subscriber is.

\paragraph{HLR, AuC and EIR} The HLR stores the permanent subscription profile and the
address of the VLR currently serving the subscriber. The AuC holds $K_i$ and computes
authentication triplets. The EIR holds lists of white, grey and black-listed IMEIs.

\subsection{Interfaces}
\begin{table}[htbp]
\centering\small
\caption{Main GSM interfaces and their protocols \cite{ts23002}.}\label{tab:interfaces}
\rowcolors{2}{acnmist}{white}
\begin{tabularx}{\linewidth}{@{}l l l Y@{}}
\toprule
\rowcolor{white}Interface & Between & Main protocol & Role \\
\midrule
Um   & MS -- BTS        & LAPDm, RR, MM, CC      & Radio interface, TDMA/FDMA, GMSK \\
Abis & BTS -- BSC       & LAPD over E1/T1        & Transport of traffic and BTS control \\
A    & BSC -- MSC       & BSSAP over SS7         & Call and mobility signalling, 64\,kbit/s circuits \\
B    & MSC -- VLR       & MAP                    & Internal in practice (MSC/VLR combined) \\
C    & GMSC -- HLR      & MAP                    & Routing interrogation for incoming calls \\
D    & VLR -- HLR       & MAP                    & Location updating, subscriber data \\
E    & MSC -- MSC       & MAP + ISUP             & Inter-MSC handover, call routing \\
F    & MSC -- EIR       & MAP                    & IMEI check \\
\bottomrule
\end{tabularx}
\end{table}

The value of an open interface is economic as well as technical: because the A
interface is fully standardised, an operator can connect any vendor's BSS to any
vendor's MSC. The Abis interface was standardised too, but in practice it often
contains vendor-specific operation and maintenance extensions.

\subsection{The radio interface Um}
GSM uses Gaussian Minimum Shift Keying (GMSK) with a bandwidth--time product of
$BT = 0.3$. GMSK is a constant-envelope modulation, so the handset power amplifier
can run near saturation, which saves battery. One bit is carried per symbol, at a
gross rate of $13/48\,\mathrm{MHz} \approx \qty{270.833}{\kilo\bit\per\second}$. The
time structure derived from this rate is drawn in \cref{fig:tdma} and computed in
\cref{sec:models}.

A normal burst carries two blocks of 57 data bits, a 26-bit training sequence in the
middle (the receiver uses it to estimate the multipath channel for its equaliser),
two stealing flags, tail bits and a guard period of 8.25 bit durations that absorbs
timing uncertainty and power ramping. Frequency hopping (up to 217 hops per second,
one hop per TDMA frame) averages fading and interference across carriers.

\subsection{Mobility: location areas and the TMSI}
Cells are grouped into \emph{location areas} (LAs). An idle MS listens to the
broadcast channel and performs a \emph{location update} only when it enters a new
LA, or periodically. When a call arrives, the network pages the MS in all cells of
its current LA. The size of an LA is a trade-off between update signalling (small
LAs) and paging load (large LAs); this trade-off is quantified in Lectures~4 and~12.

\subsection{Security: authentication and ciphering}
GSM authentication is a \emph{challenge--response} protocol \cite{ts43020}. The AuC
draws a 128-bit random number RAND and computes
\begin{subequations}\label{eq:a3a8}
\begin{align}
  \text{SRES} &= \mathrm{A3}(K_i,\text{RAND}), \qquad &&\text{(32 bits)} \label{eq:a3}\\
  K_c &= \mathrm{A8}(K_i,\text{RAND}), \qquad &&\text{(64 bits)} \label{eq:a8}
\end{align}
\end{subequations}
and sends the triplet $(\text{RAND},\text{SRES},K_c)$ to the VLR. The VLR sends only
RAND to the MS; the SIM repeats the computation of \cref{eq:a3} and returns its
result. If it matches SRES, the subscriber is authentic and both sides now share
$K_c$, which keys the A5 stream cipher on the radio link. Note what this design
achieves and what it does not:
\begin{itemize}
  \item $K_i$ never travels over any interface, and RAND is useless without $K_i$.
  \item Authentication is \emph{one-way}: the network is not authenticated to the MS,
        so a false base station can impersonate the network.
  \item Encryption ends at the BTS; Abis, A and the core links are not protected by A5.
  \item $K_c$ has only 64 bits, and the A5/1 and A5/2 ciphers have published
        cryptanalytic attacks; A5/2 has been withdrawn and A5/3 (based on KASUMI)
        was added later.
\end{itemize}
UMTS (Lecture~3) and LTE (Lectures~9 and~11) fix these weaknesses with mutual
authentication, integrity protection and longer keys.

\subsection{GPRS: a packet overlay}
GPRS adds packet switching while reusing the GSM radio carriers, sites and
subscriber databases \cite{ts23060}. Three new functions are needed:
\begin{itemize}
  \item \textbf{Packet Control Unit (PCU)} in the BSS, which implements the
        packet-oriented RLC/MAC layer and schedules radio blocks among users.
  \item \textbf{Serving GPRS Support Node (SGSN)}, the packet counterpart of the MSC:
        mobility management (routing areas), authentication, ciphering (in GPRS,
        ciphering runs between MS and SGSN at the LLC layer) and charging.
  \item \textbf{Gateway GPRS Support Node (GGSN)}, the border router to external
        packet data networks: it allocates the user IP address and anchors the session.
\end{itemize}
The SGSN connects to the BSS over \iface{Gb} (originally Frame Relay), to the GGSN
over \iface{Gn} using the GPRS Tunnelling Protocol (GTP), and to the HLR over
\iface{Gr} using MAP. The GGSN connects to the Internet or a corporate network over
\iface{Gi}. The idea of a \emph{tunnel} between a serving node and a gateway node is
the single most durable concept introduced by GPRS: it survives as the S5/S8
interface of LTE and the N9 interface of 5G.

On the radio side, GPRS defines packet data channels (PDCH). A PDCH is an ordinary
timeslot dedicated to packet traffic; data is sent in \emph{radio blocks} of four
consecutive bursts (one block every \qty{20}{\milli\second} on average). Four coding
schemes trade protection for rate:
CS-1 (\qty{9.05}{\kilo\bit\per\second}), CS-2 (13.4), CS-3 (15.6) and CS-4
(\qty{21.4}{\kilo\bit\per\second}) per timeslot. A multislot-capable MS can use
several timeslots at once.

\subsection{GPRS mobility management states}
A GPRS MS is in one of three mobility-management states:
\textbf{IDLE} (not attached; the network does not know the MS),
\textbf{STANDBY} (attached; location known per routing area; must be paged before
data can be delivered) and \textbf{READY} (location known per cell; data can be sent
immediately). A timer moves the MS from READY back to STANDBY after a period of
inactivity. The same idea --- keep the device reachable cheaply while idle, and
spend signalling only when traffic flows --- reappears as ECM-IDLE/CONNECTED in LTE and
CM-IDLE/CONNECTED in 5G.

\subsection{EDGE: higher-order modulation on the same carrier}
EDGE keeps the GSM symbol rate ($270.833$\,ksymbol/s), the \qty{200}{\kilo\hertz}
carrier and the TDMA structure, but introduces 8-PSK, which carries three bits per
symbol instead of one. Enhanced GPRS (EGPRS) defines nine modulation and coding
schemes: MCS-1 to MCS-4 use GMSK and MCS-5 to MCS-9 use 8-PSK, from
\qty{8.8}{\kilo\bit\per\second} up to \qty{59.2}{\kilo\bit\per\second} per timeslot.
Two mechanisms make the higher rates usable:
\begin{itemize}
  \item \textbf{Link adaptation:} the network selects the MCS from channel-quality
        measurements, moving to a robust scheme when the user is at the cell edge.
  \item \textbf{Incremental redundancy:} if a block is received in error, the
        retransmission carries different puncturing of the same coded bits, and
        the receiver combines both copies. This is an early form of the hybrid ARQ
        used in HSPA, LTE and NR.
\end{itemize}

% ---------------------------------------------------------------------
\section{Mathematical models}\label{sec:models}

\subsection{GSM radio timing}\label{sec:timing}
\paragraph{Engineering problem} Every timing parameter of GSM --- the burst length,
the frame duration, the maximum cell radius --- follows from one number, the gross
bit rate. We derive them to see where the familiar values come from.

\paragraph{System model and assumptions} One carrier, GMSK with one bit per symbol,
gross rate $R = f_{\mathrm{ref}}/48$ with reference clock $f_{\mathrm{ref}} = \qty{13}{\mega\hertz}$.
A timeslot lasts 156.25 bit periods and a TDMA frame contains eight timeslots.

\paragraph{Derivation} The bit period and the derived durations are
\begin{align}
  T_{\mathrm{b}} &= \frac{1}{R} = \frac{48}{13\times10^{6}}\ \mathrm{s}
     \approx \qty{3.692}{\micro\second}, \label{eq:tb}\\
  T_{\mathrm{slot}} &= 156.25\,T_{\mathrm{b}} = \frac{15}{26}\ \mathrm{ms}
     \approx \qty{576.9}{\micro\second}, \label{eq:tslot}\\
  T_{\mathrm{frame}} &= 8\,T_{\mathrm{slot}} = \frac{60}{13}\ \mathrm{ms}
     \approx \qty{4.615}{\milli\second}, \label{eq:tframe}\\
  T_{26} &= 26\,T_{\mathrm{frame}} = \qty{120}{\milli\second}. \label{eq:t26}
\end{align}
\paragraph{Gross rate of a full-rate traffic channel} In each 26-multiframe a
full-rate TCH uses 24 frames for traffic, one for the SACCH and leaves one idle.
Each traffic burst carries $2\times 57 = 114$ coded bits, so
\begin{equation}\label{eq:tchrate}
  R_{\mathrm{TCH/F}} = \frac{114\ \mathrm{bit}\times 24}{T_{26}}
  = \frac{2736\ \mathrm{bit}}{0.120\ \mathrm{s}} = \qty{22.8}{\kilo\bit\per\second}.
\end{equation}
The 13\,kbit/s speech codec output is protected by channel coding that expands it to
exactly this 22.8\,kbit/s.

\paragraph{Physical interpretation} The ratio
$R_{\mathrm{TCH/F}}/R = 22.8/270.833 \approx 8.4\%$ is close to $1/8$ minus the share of
the frame spent on training sequence, tail bits, guard time and signalling. GSM
therefore spends roughly a quarter of its gross capacity on overhead --- a figure to
remember when comparing it with OFDM systems in Lectures~6 and~14.

\paragraph{Spectral efficiency} With $B = \qty{200}{\kilo\hertz}$ the gross
spectral efficiency is $\eta = R/B \approx \qty{1.35}{bit/s/Hz}$. The Shannon bound
\cite{shannon1948mathematical}
\begin{equation}\label{eq:shannon}
  C = B\log_2\!\left(1+\SNR\right)
\end{equation}
with $C$ the capacity in bit/s, $B$ the bandwidth in Hz and $\SNR$ the linear
signal-to-noise ratio, shows that $1.35$\,bit/s/Hz could in principle be carried at
$\SNR = 2^{1.354}-1 \approx 1.56$ ($\approx\qty{1.9}{\dB}$). GSM operates at a far
higher SNR (around \qty{9}{\dB} carrier-to-interference) because it uses simple,
robust modulation and coding rather than capacity-approaching codes. \Cref{eq:shannon}
assumes an ideal band-limited additive white Gaussian noise (AWGN) channel and
unlimited coding complexity; it is an upper bound, not a prediction of throughput.

\subsection{Timing advance and the maximum cell radius}
\paragraph{Problem} The BTS expects each burst to arrive inside its timeslot. A
distant MS sees the downlink late and its uplink arrives late again, so it must
transmit early by the round-trip delay.
\paragraph{Model} Line-of-sight propagation at speed $c$; the timing advance (TA) is
signalled in steps of one bit period, $\mathrm{TA}\in\{0,1,\dots,63\}$.
\paragraph{Derivation} A round-trip delay of $2d/c$ must be compensated by
$\mathrm{TA}\cdot T_{\mathrm{b}}$, so one TA step corresponds to a distance
\begin{equation}\label{eq:tastep}
  \Delta d = \frac{c\,T_{\mathrm{b}}}{2} = \frac{\qty{2.998e8}{\metre\per\second}\times\qty{3.692}{\micro\second}}{2}
  \approx \qty{553}{\metre},
  \qquad d_{\max} = 63\,\Delta d \approx \qty{34.9}{\kilo\metre}.
\end{equation}
\paragraph{Interpretation and limits} The familiar ``35\,km'' limit of a GSM cell is a
consequence of the 6-bit TA field, not of the link budget. Extended-range cells
overcome it by using two adjacent timeslots per user, at the cost of capacity. The
same reasoning, with a finer step, gives the LTE timing advance in Lecture~8.

\subsection{Co-channel interference and the cluster size}\label{sec:reuse}
\paragraph{Problem} How tightly can carriers be reused? Smaller reuse distance means
more carriers per cell but stronger interference.
\paragraph{Model and assumptions} Hexagonal cells of radius $R_{\mathrm{c}}$, clusters
of $N = i^2 + ij + j^2$ cells ($i,j$ non-negative integers), power-law path loss
$\Prx \propto d^{-n}$ with exponent $n$, a user at the cell edge, and only the six
first-tier co-channel cells at reuse distance $D = R_{\mathrm{c}}\sqrt{3N}$ \cite{rappaport2002wireless}.
\paragraph{Derivation} The wanted signal is $C \propto R_{\mathrm{c}}^{-n}$ and each of the
six interferers contributes approximately $D^{-n}$, hence
\begin{equation}\label{eq:ci}
  \frac{C}{I} \approx \frac{R_{\mathrm{c}}^{-n}}{6\,D^{-n}}
  = \frac{1}{6}\left(\frac{D}{R_{\mathrm{c}}}\right)^{n}
  = \frac{\left(\sqrt{3N}\right)^{n}}{6}.
\end{equation}
\paragraph{Interpretation} \Cref{eq:ci} grows with both $N$ and $n$: a large cluster
puts interferers further away, and a steep path loss attenuates them faster. Since
GSM tolerates a $C/I$ of roughly \qty{9}{\dB}, \cref{fig:reuse} suggests that $N=4$
($13.8$\,dB for $n=4$) is comfortable, and $N=3$ is feasible for $n=4$ but marginal
for $n=3.5$. In practice operators combined tight reuse with frequency hopping and
power control, which average and reduce interference.
\paragraph{Limitations} The model ignores shadow fading (which varies $C/I$ by several dB),
sectorisation (which removes some interferers), and the fact that not all
interferers are active at once. It is a planning first estimate, not a guarantee.

\begin{acnfigure}
  \acnfig{lecture-01/reuse-ci}
  \caption{First-tier $C/I$ from \cref{eq:ci} versus cluster size for two path-loss
  exponents. The horizontal line marks the approximate GSM protection ratio.}
  \label{fig:reuse}
\end{acnfigure}

\subsection{Packet throughput with multiple timeslots}
For a user allocated $n_{\mathrm{TS}}$ timeslots with per-timeslot rate $r_{\mathrm{CS}}$,
the peak radio-link-control (RLC) data rate is
\begin{equation}\label{eq:gprsrate}
  R_{\mathrm{user}} = n_{\mathrm{TS}}\; r_{\mathrm{CS}}\;(1-p_{\mathrm{BLER}}),
\end{equation}
where $p_{\mathrm{BLER}}$ is the block error ratio and the factor $(1-p_{\mathrm{BLER}})$ assumes
every erroneous block is retransmitted once successfully (a simple ARQ model). Higher
coding schemes increase $r_{\mathrm{CS}}$ but, at a given $C/I$, also increase
$p_{\mathrm{BLER}}$; the best scheme maximises the product. \Cref{fig:rates} shows the
peak values.

\begin{acnfigure}
  \acnfig{lecture-01/gprs-rates}
  \caption{Peak per-user RLC data rate versus number of timeslots ($p_{\mathrm{BLER}}=0$)
  for two GPRS and two EGPRS schemes. Real throughput is lower because of
  retransmissions, sharing among users and higher-layer overhead.}
  \label{fig:rates}
\end{acnfigure}

% ---------------------------------------------------------------------
\section{Network architecture and diagrams}\label{sec:architecture}
\Cref{fig:arch} combines the circuit-switched GSM network and the GPRS packet
overlay. Two observations are worth making. First, the BSS is shared: the same BTS
carries voice timeslots towards the MSC and packet timeslots towards the SGSN, and
the BSC splits them. Second, the HLR is shared: a single subscription record serves
both domains, which is why GPRS could be launched without re-provisioning subscribers.

\begin{acnfigure}
  \acnfig{lecture-01/gsm-architecture}
  \caption{GSM network with the GPRS overlay. Solid grey: transport links;
  teal: packet-switched links; dashed blue: MAP signalling towards the HLR and EIR.}
  \label{fig:arch}
\end{acnfigure}

\begin{acnfigure}
  \resizebox{\linewidth}{!}{\acnfig{lecture-01/tdma-hierarchy}}
  \caption{GSM time structure. A 26-multiframe (\qty{120}{\milli\second}) contains 24
  traffic frames, one SACCH frame and one idle frame; each TDMA frame contains eight
  timeslots; a normal burst carries 114 data bits.}
  \label{fig:tdma}
\end{acnfigure}

% ---------------------------------------------------------------------
\section{Step-by-step procedures}\label{sec:procedures}

\begin{procedurebox}{Authentication and start of ciphering}\label{proc:auth}
\begin{steps}
  \item The MS sends a \msg{Location Updating Request} (or call set-up request)
        identifying itself with its TMSI, or with its IMSI if no TMSI is valid.
  \item If the VLR has no unused triplet, it requests a batch from the HLR/AuC with MAP
        \msg{Send Authentication Info}; the AuC returns several triplets.
  \item The MSC/VLR sends \msg{Authentication Request} containing only RAND.
  \item The SIM computes SRES$'$ and $K_c$ from \cref{eq:a3a8} and the MS returns SRES$'$.
  \item The VLR compares SRES$'$ with SRES. On a match it orders ciphering by sending
        $K_c$ to the BSS in a BSSMAP \msg{Cipher Mode Command}.
  \item The BTS sends \msg{Ciphering Mode Command}; the MS starts A5 with $K_c$ and
        answers \msg{Ciphering Mode Complete}. From now on Um is encrypted.
  \item The VLR may allocate a new TMSI, sent encrypted, so the IMSI is protected next time.
\end{steps}
\end{procedurebox}

\begin{acnfigure}
  \resizebox{\linewidth}{!}{\acnfig{lecture-01/gsm-authentication}}
  \caption{Message sequence of GSM authentication and ciphering (\cref{proc:auth}).}
  \label{fig:auth}
\end{acnfigure}

\begin{procedurebox}{GPRS attach and PDP context activation}\label{proc:pdp}
\begin{steps}
  \item \textbf{GPRS attach.} The MS sends \msg{Attach Request} to the SGSN. The SGSN
        authenticates the MS (triplets from the HLR over Gr), registers itself in the
        HLR with \msg{Update GPRS Location}, receives the subscription with
        \msg{Insert Subscriber Data} and returns \msg{Attach Accept} with a P-TMSI.
        The MS is now in READY state but has no IP address.
  \item \textbf{Session request.} The MS sends \msg{Activate PDP Context Request}
        with an NSAPI (identifier of the context), the requested access point name
        (APN) and QoS.
  \item \textbf{Gateway selection.} The SGSN checks the subscription and resolves the
        APN with DNS to the address of a GGSN.
  \item \textbf{Tunnel set-up.} The SGSN sends GTP-C \msg{Create PDP Context Request};
        the GGSN allocates an IP address (or obtains one from a RADIUS/DHCP server)
        and returns its tunnel endpoint identifier (TEID).
  \item \textbf{Completion.} The SGSN sends \msg{Activate PDP Context Accept} with the
        IP address. User packets now flow: MS--SGSN over LLC/SNDCP, SGSN--GGSN inside a
        GTP-U tunnel identified by TEIDs.
\end{steps}
\end{procedurebox}

\begin{acnfigure}
  \resizebox{\linewidth}{!}{\acnfig{lecture-01/pdp-activation}}
  \caption{GPRS attach followed by PDP context activation (\cref{proc:pdp}).
  Dashed blue: signalling; solid teal: user plane.}
  \label{fig:pdp}
\end{acnfigure}

The procedure illustrates a principle that the course will meet again in LTE (the
default bearer, Lecture~9) and 5G (the PDU session, Lecture~20): \emph{being
attached} (known to the network) and \emph{having connectivity} (an IP address and a
tunnel) are separate states, established by separate procedures. LTE later merged
them for the first PDN connection to reduce delay.

% ---------------------------------------------------------------------
\section{Worked examples}\label{sec:examples}

\begin{examplebox}{Capacity of a GSM-900 operator}\label{ex:capacity}
\textbf{Problem.} An operator owns \qty{10}{\mega\hertz} of paired GSM-900 spectrum and
plans with cluster size $N=4$. How many full-rate voice channels does a cell offer?

\textbf{Solution.} Number of carriers: $\floor{\qty{10}{\mega\hertz}/\qty{200}{\kilo\hertz}} = 50$;
one is kept as a guard, so 49 are usable. Per cell: $\floor{49/4} = 12$ carriers,
i.e.\ $12\times 8 = 96$ timeslots. Reserving one timeslot for the broadcast/common
control channel and one for SDCCH signalling leaves about $94$ traffic channels.

\textbf{Interpretation.} 94 simultaneous calls per cell is the hard limit of circuit
switching. How much \emph{traffic} (in erlangs) these channels carry at a given
blocking probability is computed with the Erlang-B formula in Lecture~2.
\end{examplebox}

\begin{examplebox}{Timing advance}\label{ex:ta}
\textbf{Problem.} The BTS measures that an MS's bursts arrive \qty{55}{\micro\second}
late. Which TA value is signalled and how far is the MS?

\textbf{Solution.} The round trip lasts $\qty{55}{\micro\second}$, so
$\mathrm{TA} = \mathrm{round}(55/3.692) = 15$ and
$d \approx 15\times\qty{553}{\metre} \approx \qty{8.3}{\kilo\metre}$.
(Directly: $d = c\,\tau/2 = 3\times10^8\times55\times10^{-6}/2 = \qty{8.25}{\kilo\metre}$.)

\textbf{Interpretation.} The TA quantisation error is up to half a step, $\approx\qty{277}{\metre}$;
TA therefore also gives the network a coarse estimate of the user's distance, the
basis of early cell-ID+TA positioning.
\end{examplebox}

\begin{examplebox}{Choosing a GPRS coding scheme}\label{ex:cs}
\textbf{Problem.} At a given $C/I$ a user observes $p_{\mathrm{BLER}} = 0.05$ with CS-2
and $p_{\mathrm{BLER}} = 0.45$ with CS-4. With four timeslots, which scheme gives the
higher throughput according to \cref{eq:gprsrate}?

\textbf{Solution.}
\begin{align*}
  R_{\mathrm{CS\text{-}2}} &= 4\times13.4\times(1-0.05) = \qty{50.9}{\kilo\bit\per\second},\\
  R_{\mathrm{CS\text{-}4}} &= 4\times21.4\times(1-0.45) = \qty{47.1}{\kilo\bit\per\second}.
\end{align*}
\textbf{Interpretation.} CS-2 wins despite its lower nominal rate. Link adaptation is
precisely the selection of the scheme that maximises this product; the same logic
reappears with CQI and MCS tables in LTE (Lecture~7).
\end{examplebox}

% ---------------------------------------------------------------------
\section{Technology comparison}\label{sec:comparison}
\begin{table}[htbp]
\centering\small
\caption{GSM, GPRS and EDGE compared. Rates are peak RLC rates per timeslot or per
user with eight timeslots; achievable user rates are lower.}\label{tab:compare}
\rowcolors{2}{acnmist}{white}
\begin{tabularx}{\linewidth}{@{}l Y Y Y@{}}
\toprule
\rowcolor{white} & GSM (circuit data) & GPRS & EDGE (EGPRS) \\
\midrule
Switching          & Circuit                 & Packet                     & Packet \\
Modulation         & GMSK                    & GMSK                       & GMSK and 8-PSK \\
Bits per symbol    & 1                       & 1                          & 1 or 3 \\
Rate per timeslot  & 9.6 or 14.4\,kbit/s     & 9.05--21.4\,kbit/s         & 8.8--59.2\,kbit/s \\
Peak (8 TS)        & --                      & 171.2\,kbit/s              & 473.6\,kbit/s \\
Resource use       & Reserved for the call   & Shared, on demand          & Shared, on demand \\
Retransmission     & None (FEC only)         & ARQ in RLC                 & ARQ + incremental redundancy \\
Core nodes         & MSC, GMSC               & SGSN, GGSN                 & SGSN, GGSN (unchanged) \\
Charging basis     & Time                    & Volume                     & Volume \\
\bottomrule
\end{tabularx}
\end{table}

\paragraph{Trade-offs} GPRS made data cheap (volume-based charging, no idle circuit)
but introduced queuing delay and variable throughput; round-trip times of several
hundred milliseconds were typical. EDGE tripled the peak rate with a software and
transceiver upgrade, but 8-PSK is not constant-envelope, so it needs more linear power
amplifiers, and its high schemes are only usable close to the BTS. Neither changed
the fundamental limit: data rides on 200\,kHz carriers designed for voice. Breaking
that limit required a new radio interface, WCDMA in 3G (Lecture~3).

% ---------------------------------------------------------------------
\section{Key takeaways}\label{sec:takeaways}
\begin{takeawaybox}
\begin{itemize}
  \item GSM is a \emph{system} standard: open interfaces (A, Um) and a separable
        identity (SIM/IMSI) enabled multi-vendor networks and roaming.
  \item All GSM timing follows from $R = \qty{270.833}{\kilo\bit\per\second}$: bit
        $\qty{3.69}{\micro\second}$, slot $\qty{577}{\micro\second}$, frame $\qty{4.615}{\milli\second}$,
        26-multiframe $\qty{120}{\milli\second}$; a full-rate TCH carries 22.8\,kbit/s gross.
  \item The 6-bit timing advance limits a normal cell to about \qty{35}{\kilo\metre}.
  \item Authentication uses triplets $(\text{RAND},\text{SRES},K_c)$; it is one-way, and
        encryption stops at the BTS.
  \item GPRS adds PCU, SGSN and GGSN; the PDP context and the GTP tunnel are the
        ancestors of the LTE bearer and the 5G PDU session.
  \item EDGE raises the bit rate per symbol (8-PSK) and adds incremental redundancy,
        an early hybrid ARQ.
  \item The first-tier model $C/I \approx (\sqrt{3N})^n/6$ links reuse, path loss and
        interference; it is a first planning estimate.
\end{itemize}
\end{takeawaybox}

% ---------------------------------------------------------------------
\section{Exercises}\label{sec:exercises}

\begin{exercisebox}{1}{Architecture}\label{exr:arch}
Explain why the HLR stores the address of the VLR but not the cell in which the
subscriber is located. What would be the signalling consequence of storing the cell?
\end{exercisebox}

\begin{exercisebox}{1}{Security}\label{exr:sec}
List two attacks made possible by the one-way authentication of GSM and say which
later mechanism (Lecture~11) prevents each.
\end{exercisebox}

\begin{exercisebox}{2}{Half-rate channel}\label{exr:halfrate}
A half-rate traffic channel uses a timeslot in every second traffic frame. Using
\cref{eq:tchrate}, compute its gross bit rate, and explain how it doubles the voice
capacity of the cell in \cref{ex:capacity}.
\end{exercisebox}

\begin{exercisebox}{2}{Reuse planning}\label{exr:reuse}
Using \cref{eq:ci}, find the smallest valid cluster size $N$ that guarantees
$C/I\ge\qty{12}{\dB}$ for (a) $n=4$ and (b) $n=3$.
\end{exercisebox}

\begin{exercisebox}{2}{EDGE throughput}\label{exr:edge}
A user has five timeslots. With MCS-9 the block error ratio is 0.30; with MCS-7
(\qty{44.8}{\kilo\bit\per\second} per timeslot) it is 0.08. Which scheme should link
adaptation choose, and what is the expected throughput?
\end{exercisebox}

\begin{exercisebox}{3}{Overhead budget}\label{exr:overhead}
Starting from \cref{eq:tb,eq:tframe}, compute what fraction of the gross bit rate of
one timeslot is used by (i) guard period, (ii) tail bits and stealing flags, (iii)
training sequence, (iv) SACCH and idle frames. Compare the total with the cyclic-prefix
overhead of LTE ($\approx 7\%$) and discuss why the comparison is not entirely fair.
\end{exercisebox}

\begin{exercisebox}{3}{Generalising the C/I model}\label{exr:ci3}
Extend \cref{eq:ci} to a cell divided into three \ang{120} sectors, where only two
first-tier interferers fall into the main lobe of the sector antenna. Derive the new
expression and the gain in dB, and discuss which assumptions of the original model
the result inherits.
\end{exercisebox}

\subsection*{Solutions}
\addcontentsline{toc}{subsection}{Solutions}

\begin{solutionbox}{\ref{exr:arch}}
The cell changes every few minutes for a moving user; if the HLR tracked it, every
cell change of every subscriber in the world would generate signalling to the home
network, often internationally. Storing only the VLR address limits HLR updates to
changes of MSC area, and the location area kept in the VLR limits radio paging to a
group of cells. This two-level hierarchy (HLR $\to$ VLR $\to$ LA) is the core of GSM
mobility management.
\end{solutionbox}

\begin{solutionbox}{\ref{exr:sec}}
(1) A false base station can request the IMSI in clear and track the user (IMSI
catching); 5G conceals the permanent identity with SUCI, and LTE/UMTS reduce the
exposure. (2) A false base station can accept the MS without authentication and
order no ciphering or a weak cipher, then relay traffic (man-in-the-middle). UMTS/LTE
mutual authentication with the AUTN token (Lecture~11) lets the USIM verify that the
challenge was produced by its home network, and integrity protection of signalling
prevents tampering with the security-mode command.
\end{solutionbox}

\begin{solutionbox}{\ref{exr:halfrate}}
A half-rate channel gets 12 of the 24 traffic frames per 26-multiframe, so
$R_{\mathrm{TCH/H}} = 114\times 12/0.120 = \qty{11.4}{\kilo\bit\per\second}$. Two half-rate
users share one timeslot, so the 94 timeslots of \cref{ex:capacity} can carry up to
188 calls, at the price of a lower-rate (lower-quality) speech codec.
\end{solutionbox}

\begin{solutionbox}{\ref{exr:reuse}}
We need $(\sqrt{3N})^n/6 \ge 10^{1.2} = 15.85$, i.e.\ $(3N)^{n/2}\ge 95.1$.
(a) $n=4$: $(3N)^2\ge 95.1 \Rightarrow N\ge 3.25$; the smallest valid $N$ is 4
($C/I=\qty{13.8}{\dB}$). (b) $n=3$: $(3N)^{1.5}\ge95.1 \Rightarrow 3N\ge 20.8 \Rightarrow N\ge 6.94$;
the smallest valid $N$ is 7 ($C/I = 10\log_{10}(21^{1.5}/6) = \qty{12.05}{\dB}$).
Valid sizes are $N = i^2+ij+j^2 \in\{1,3,4,7,9,12,\dots\}$.
\end{solutionbox}

\begin{solutionbox}{\ref{exr:edge}}
$R_{\mathrm{MCS\text{-}9}} = 5\times59.2\times0.70 = \qty{207.2}{\kilo\bit\per\second}$ and
$R_{\mathrm{MCS\text{-}7}} = 5\times44.8\times0.92 = \qty{206.1}{\kilo\bit\per\second}$.
The two are practically equal; MCS-9 is marginally better on throughput, but MCS-7
gives fewer retransmissions and therefore lower delay and jitter. With incremental
redundancy the effective MCS-9 throughput would be higher than the simple model
predicts, because failed blocks still contribute soft information.
\end{solutionbox}

\begin{solutionbox}{\ref{exr:overhead}}
Per 156.25-bit slot: guard $8.25/156.25 = 5.3\%$; tail and flags $(6+2)/156.25 = 5.1\%$;
training $26/156.25 = 16.6\%$; data $114/156.25 = 73.0\%$. Of the 26 frames, 2 do not
carry traffic, so the useful fraction is $0.730\times24/26 = 67.4\%$:
$270.833/8\times0.674 = \qty{22.8}{\kilo\bit\per\second}$, consistent with \cref{eq:tchrate}.
Total overhead $\approx 33\%$. LTE's cyclic prefix costs about $7\%$, but LTE also
spends resources on reference signals, control region and guard bands; a fair
comparison must count all of them (Lecture~7).
\end{solutionbox}

\begin{solutionbox}{\ref{exr:ci3}}
With two interferers, $C/I \approx (\sqrt{3N})^n/2$, which is three times (\qty{4.8}{\dB})
better than \cref{eq:ci}. For $N=4$, $n=4$: $10\log_{10}(144/2) = \qty{18.6}{\dB}$. The
result still assumes ideal hexagons, equal powers, no shadowing and a worst-case
user at the cell edge; in addition it assumes ideal antenna patterns with no
back-lobe leakage. Sectorisation trades trunking efficiency (each sector has fewer
channels) for interference reduction.
\end{solutionbox}

% ---------------------------------------------------------------------
\section{Further reading}\label{sec:reading}
The architecture of all 3GPP systems is summarised in \cite{ts23002}; the GSM radio
structure is specified in \cite{ts45002} and security functions in \cite{ts43020}.
GPRS stage-2 procedures are in \cite{ts23060} and layer-3 messages in \cite{ts24008}.
The textbook by Eberspächer et al.\ \cite{eberspaecher2009gsm} remains the most
complete single description of GSM and GPRS, and Rappaport \cite{rappaport2002wireless}
derives the cellular reuse model used in \cref{sec:reuse}. MAP is specified in
\cite{ts29002}; identities in \cite{ts23003}.

\nocite{shannon1948mathematical}
\printbibliography[heading=subbibliography,title={\ACNstrRefs}]

\end{document}
